\bichapter{RedHawk 与 RedHawk-SC Fusion}{RedHawk and RedHawk-SC Fusion}
\label{chap:redhawk}

RedHawk\texttrademark{} 与 RedHawk-SC\texttrademark{} 电源完整性（power integrity）解决方案通过 RedHawk Fusion 与 RedHawk-SC Fusion 接口集成到实现流程中。可在电源规划（power planning）与初始布局（initial placement）完成后，在物理实现流程的不同节点使用 RedHawk Fusion 与 RedHawk-SC Fusion 进行轨分析（rail analysis）。从而在详细布线（detail routing）之前发现潜在电源网络问题，显著缩短设计周期周转时间（turnaround time）。

布局完成且 PG mesh 可用后，使用 RedHawk Fusion 或 RedHawk-SC Fusion 对电源与地网络执行压降分析（voltage drop analysis），以计算功耗并检查压降违例。也可在设计流程的其他阶段（例如详细布线之后）执行压降分析。

芯片收尾（chip finishing）完成后，使用 RedHawk Fusion 或 RedHawk-SC 执行 PG 电迁移分析（PG electromigration analysis），以检查电流密度违例。

本章说明如何在环境中使用 RedHawk Fusion 或 RedHawk-SC Fusion 执行轨分析。关于 RedHawk 或 RedHawk-SC 分析流程与命令的详细说明，请参阅 \textit{RedHawk User Manual} 或 \textit{RedHawk-SC User Manual}。

本章包含以下主题：
\begin{itemize}
  \item 使用 RedHawk-SC Fusion 运行轨分析（Running Rail Analysis Using RedHawk-SC Fusion）
  \item RedHawk Fusion 与 RedHawk-SC Fusion 概述（An Overview for RedHawk Fusion and RedHawk-SC Fusion）
  \item 设置可执行程序（Setting Up the Executables）
  \item 指定 RedHawk 与 RedHawk-SC 工作目录（Specifying RedHawk and RedHawk-SC Working Directories）
  \item 为轨分析准备设计与输入数据（Preparing Design and Input Data for Rail Analysis）
  \item 将理想电压源指定为 Tap（Specifying Ideal Voltage Sources as Taps）
  \item 缺失 Via 与未连接 Pin 检查（Missing Via and Unconnected Pin Checking）
  \item 使用多个轨场景运行轨分析（Running Rail Analysis with Multiple Rail Scenarios）
  \item 执行压降分析（Performing Voltage Drop Analysis）
  \item 执行 PG 电迁移分析（Performing PG Electromigration Analysis）
  \item 执行最小路径电阻分析（Performing Minimum Path Resistance Analysis）
  \item 执行有效电阻分析（Performing Effective Resistance Analysis）
  \item 执行分布式 RedHawk Fusion 轨分析（Performing Distributed RedHawk Fusion Rail Analysis）
  \item 电压驱动的电源开关单元尺寸调整（Voltage Driven Power Switch Cell Sizing）
  \item 使用宏模型（Working With Macro Models）
  \item 执行签核分析（Performing Signoff Analysis）
  \item 写出分析与检查报告（Writing Analysis and Checking Reports）
  \item 在 GUI 中显示地图（Displaying Maps in the GUI）
  \item 在 GUI 中显示 ECO 形状（Displaying ECO Shapes in the GUI）
  \item 电压热点分析（Voltage Hotspot Analysis）
  \item 查询属性（Querying Attributes）
\end{itemize}

% ============================================================
\section{使用 RedHawk-SC Fusion 运行轨分析}
\subsection*{Running Rail Analysis Using RedHawk-SC Fusion}

与 RedHawk Fusion 能力类似，RedHawk-SC Fusion 在 Fusion Compiler 环境中支持门级轨分析与检查能力。要启用 RedHawk-SC Fusion，需同时启用应用选项 \appopt{rail.enable\_redhawk\_sc}，并禁用应用选项 \appopt{rail.enable\_redhawk}。

关于 RedHawk-SC Fusion 的更多信息，见「设置可执行程序」与「指定 RedHawk 与 RedHawk-SC 工作目录」。

RedHawk Fusion 与 RedHawk-SC 共用同一套轨分析与检查命令，以及用于查看分析结果的同一 GUI（除若干限制外）。表~\ref{tab:rh-sc-features} 比较了在 RedHawk Fusion、RedHawk-SC Fusion 之一或二者中受支持的 \cmd{analyze\_rail} 选项。

\begin{table}[htbp]
\centering
\caption{RedHawk Fusion 与 RedHawk-SC Fusion 分析功能比较（原书 Table 65）}
\label{tab:rh-sc-features}
\small
\begin{tabularx}{\textwidth}{@{}l X X@{}}
\toprule
\textbf{选项} & \textbf{RedHawk Fusion} & \textbf{RedHawk-SC Fusion} \\
\midrule
\opt{-redhawk\_script\_file} & 支持 & 支持 \\
\opt{-voltage\_drop} & 支持 & 支持 \\
\opt{-switching\_activity} & 支持 & 支持 \\
\opt{-electromigration} & 支持 & 支持 \\
\opt{-power\_analysis} & 支持 & 支持 \\
\opt{-result\_name} & 支持 & 支持 \\
\opt{-script\_only} & 支持 & 支持 \\
\opt{-bg} & 支持 & 支持 \\
\opt{-min\_path\_resistance} & 支持 & 仅与 \opt{-voltage\_drop} 联用时支持 \\
\opt{-effective\_resistance} & 支持 & 仅与 \opt{-voltage\_drop} 联用时支持 \\
\opt{-check\_missing\_via} & 支持 & 仅与 \opt{-voltage\_drop} 联用时支持；显示缺失 via 检查结果时不遵从 IR 阈值设置 \\
\opt{-extra\_gsr\_option\_file} & 支持 & 不支持 \\
\opt{-multiple\_script\_files} & 支持 & 不支持 \\
\opt{-submit\_to\_other\_machines} & 支持 & 不支持 \\
\bottomrule
\end{tabularx}
\end{table}

在使用 RedHawk-SC Fusion 执行轨分析之前，必须按「为轨分析准备设计与输入数据」所述指定库与所需输入文件的位置。以下两个应用选项仅在 RedHawk-SC Fusion 中可用：

\begin{itemize}
  \item \appopt{rail.toggle\_rate}：为轨分析指定不同单元类型的翻转率（toggle rate）。例如：
\begin{lstlisting}
fc_shell> set_app_options -name rail.toggle_rate  \
   -value {clock 2.0 data 0.2 combinational 0.15 \
   sequential 0.15}
\end{lstlisting}
  \item \appopt{rail.em\_only\_tech\_file}：指定用于 PG 电迁移分析的电迁移规则文件。当电迁移规则信息定义在与 \appopt{rail.tech\_file} 所指定技术文件不同的单独技术文件中时，使用此选项。例如：
\begin{lstlisting}
fc_shell> set_app_options -name rail.em_only_tech_file \
   -value EM_ONLY.rule
\end{lstlisting}
\end{itemize}

% ============================================================
\section{RedHawk Fusion 与 RedHawk-SC Fusion 概述}
\subsection*{An Overview for RedHawk Fusion and RedHawk-SC Fusion}

RedHawk Fusion 与 RedHawk-SC Fusion 均支持门级轨分析能力，包括静态与动态轨分析。可在 Fusion Compiler 环境中使用二者执行下列类型的检查与分析：

\begin{itemize}
  \item \textbf{缺失 via 与未连接 pin 形状检查}\\
    要检查设计中的缺失 via 或未连接 pin 形状，先使用 \cmd{set\_missing\_via\_check\_options} 配置检查相关设置，再使用 \cmd{analyze\_rail} \opt{-check\_missing\_via} 执行检查。详见「缺失 Via 与未连接 Pin 检查」。
  \item \textbf{压降分析}\\
    使用 \cmd{analyze\_rail} \opt{-voltage\_drop}，或在 GUI 中选择 Rail $>$ Analyze Rail 并选择 ``Voltage drop analysis''。将 \opt{-voltage\_drop} 设为 \texttt{static}、\texttt{dynamic}、\texttt{dynamic\_vcd} 或 \texttt{dynamic\_vectorless} 以选择分析类型。详见「执行压降分析」。
  \item \textbf{PG 电迁移分析}\\
    使用 \cmd{analyze\_rail} \opt{-electromigration}。详见「执行 PG 电迁移分析」。
  \item \textbf{最小路径电阻分析}\\
    使用 \cmd{analyze\_rail} \opt{-min\_path\_resistance}。在 RedHawk-SC 分析流程中，必须将 \opt{-voltage\_drop} 与 \opt{-min\_path\_resistance} 一起使用。详见「执行最小路径电阻分析」。
\end{itemize}

根据所运行的分析类型，工具会生成可在 Fusion Compiler GUI 中查看的可视化显示（地图，maps），以及可在错误浏览器（error browser）中显示的错误数据。这些地图与错误数据有助于在不离开 Fusion Compiler 环境的情况下发现问题区域并确定纠正措施。

\begin{noteBox}
Fusion Compiler 工具中的 RedHawk Fusion 不支持 RedHawk 签核分析功能，例如层次化分析，或带集总（lumped）/SPICE package 的动态分析。详见「为轨分析准备设计与输入数据」。

默认情况下，分析多角多模（multicorner-multimode）设计时，RedHawk Fusion 仅分析当前设计场景（design scenario）。关于如何创建并指定用于轨分析的轨场景（rail scenario），见「使用多个轨场景运行轨分析」。
\end{noteBox}

图~\ref{fig:rh-flow} 示出典型设计流程中使用这些分析能力的位置。

\begin{figure}[htbp]
\centering
\fbox{\parbox{0.85\textwidth}{\centering\vspace{1.5cm}
\textit{（原书 Figure 198：Using RedHawk Fusion in the Fusion Compiler Design Flow）}\\[0.3em]
\small 设计规划 $\rightarrow$ 查找缺失 via；\texttt{place\_opt}/初始布局后静态压降；\texttt{clock\_opt}/filler/电源开关插入后静态或动态压降；\texttt{route\_opt} 后 PG 电迁移分析与修复。\\
请参阅原版 PDF 插图。
\vspace{1.5cm}}}
\caption{在 Fusion Compiler 设计流程中使用 RedHawk Fusion}
\label{fig:rh-flow}
\end{figure}

\subsection{RedHawk Fusion 与 RedHawk-SC Fusion 数据流}
\subsubsection*{RedHawk Fusion and RedHawk-SC Fusion Data Flow}

RedHawk/RedHawk-SC Fusion 功能允许在实现阶段执行轨分析。在具备所需输入文件的情况下，Fusion Compiler 工具创建用于调用 RedHawk 工具的 RedHawk 运行脚本与 Global System Requirements（GSR）配置文件（或用于调用 RedHawk-SC 工具的 RedHawk-SC Python 脚本），以运行 PG net 抽取、功耗分析、压降分析与 PG 电迁移分析。

分析完成后，Fusion Compiler 工具使用 RedHawk 或 RedHawk-SC 计算的结果生成分析报告与地图。随后可在 Fusion Compiler GUI 中以图形方式检查热点。还提供错误数据与 ASCII 报告，用于检查超出限值的位置。

图~\ref{fig:rh-dataflow} 说明在 Fusion Compiler 环境中使用 RedHawk Fusion 或 RedHawk-SC Fusion 执行轨分析时的数据流。

\begin{figure}[htbp]
\centering
\fbox{\parbox{0.85\textwidth}{\centering\vspace{1.5cm}
\textit{（原书 Figure 199：RedHawk Fusion and RedHawk-SC Fusion Data Flow）}\\[0.3em]
\small 请参阅原版 PDF 插图。
\vspace{1.5cm}}}
\caption{RedHawk Fusion 与 RedHawk-SC Fusion 数据流}
\label{fig:rh-dataflow}
\end{figure}

\subsection{RedHawk/RedHawk-SC Fusion 分析流程}
\subsubsection*{RedHawk/RedHawk-SC Fusion Analysis Flow}

指定必要的输入与设计数据后，即可对设计执行压降与 PG 电迁移分析。图~\ref{fig:rh-basic-steps} 说明基本静态轨分析流程中的步骤。

\begin{figure}[htbp]
\centering
\fbox{\parbox{0.85\textwidth}{\centering\vspace{1.5cm}
\textit{（原书 Figure 200：Required Steps for a Basic Rail Analysis）}\\[0.3em]
\small 读取设计数据（\texttt{open\_lib}/\texttt{open\_block}）$\rightarrow$ 设置应用选项（\texttt{set\_app\_options}）$\rightarrow$ 指定理想电压源（\texttt{create\_taps}）$\rightarrow$ 运行轨分析（\texttt{analyze\_rail -voltage\_drop}）$\rightarrow$ 检查结果；并可运行 PG 电迁移分析（\texttt{analyze\_rail -electromigration}）检查电流违例。\\
请参阅原版 PDF 插图。
\vspace{1.5cm}}}
\caption{使用 RedHawk 与 RedHawk-SC Fusion 进行基本轨分析的必要步骤}
\label{fig:rh-basic-steps}
\end{figure}

\subsection{在后台运行 RedHawk Fusion 命令}
\subsubsection*{Running RedHawk Fusion Commands in the Background}

默认情况下，执行 \cmd{analyze\_rail} 进行轨分析时，会调用 RedHawk 或 RedHawk-SC 工具以输入数据执行指定分析类型，并在完成后将分析结果加载回轨数据库。在 RedHawk Fusion 或 RedHawk-SC Fusion 进程结束前，无法运行任何 Fusion Compiler 命令。

要在工具于后台运行时执行布局编辑任务，运行 \cmd{analyze\_rail} \opt{-bg}。分析完成后，运行 \cmd{open\_rail\_result} \opt{-back\_annotate} 将分析结果上传到轨数据库。

默认情况下，RedHawk 或 RedHawk-SC 在分析完成后将结果加载回轨数据库。但若在 \cmd{analyze\_rail} 中指定了 \opt{-bg}，则需运行 \cmd{open\_rail\_result} \opt{-back\_annotate}，以根据最新分析结果重建轨数据库并创建新的 \texttt{RAIL\_DATABASE} 文件。

\begin{noteBox}
在设计中移动 net 之后回注（back-annotate）先前运行的分析结果时，基于实例的地图会反映更新后的实例位置变化，但寄生地图不会。

打开轨结果时，工具不检测该结果是否在使用 \opt{-bg} 的情况下生成。若对非 \opt{-bg} 生成的结果运行 \cmd{open\_rail\_result} \opt{-back\_annotate}（或反之），命令可能无法正确工作。
\end{noteBox}

\subsection{报告并检查后台进程状态}
\subsubsection*{Reporting and Checking the Status of the Background Process}

当 RedHawk 或 RedHawk-SC Fusion 进程结束时，工具发出如下消息：
\begin{lstlisting}
Info: Running REDHAWK_BINARY in background with log file: LOG_FILE
\end{lstlisting}
\texttt{LOG\_FILE} 保存在 PWD 目录中。

要检查后台 \cmd{analyze\_rail} 进程是否仍在活动，使用 \cmd{report\_background\_jobs} 命令。

% ============================================================
\section{设置可执行程序}
\subsection*{Setting Up the Executables}

要运行 RedHawk Fusion 或 RedHawk-SC Fusion 功能，通过设置 \appopt{rail.product} 与 \appopt{rail.redhawk\_path} 应用选项启用功能并指定 RedHawk 或 RedHawk-SC 可执行程序的位置。

例如：
\begin{lstlisting}
fc_shell> set_app_options -name rail.product -value redhawk|redhawk_sc
fc_shell> set_app_options -name rail.redhawk_path \
    -value /tools/RedHawk_Linux64e5_V19.0.2p2/bin
\end{lstlisting}

必须确保所指定的可执行程序与所用 Fusion Compiler 版本兼容。

\begin{seeAlsoBox}
指定 RedHawk 与 RedHawk-SC 工作目录
\end{seeAlsoBox}

% ============================================================
\section{指定 RedHawk 与 RedHawk-SC 工作目录}
\subsection*{Specifying RedHawk and RedHawk-SC Working Directories}

轨分析期间，RedHawk/RedHawk-SC Fusion 创建工作目录以存放生成的文件，包括分析日志、脚本与各类输出数据。默认情况下，RedHawk 或 RedHawk-SC 工作目录名为 \texttt{RAIL\_DATABASE}。要使用不同名称，使用 \appopt{rail.database} 应用选项。

数据结构如下：
\begin{itemize}
  \item \texttt{RAIL\_CHECKING\_DIR}：RedHawk Fusion 保存轨分析期间缺失信息相关报告文件的目录。例如，运行静态轨分析时，工具在该目录中生成：
    \begin{itemize}
      \item \texttt{libcell.apl\_current}
      \item \texttt{libcell.apl\_cap}
      \item \texttt{libcell.missing\_liberty}
    \end{itemize}
  \item \texttt{RAIL\_DATABASE}：工具保存并检索功耗计算与轨分析期间生成的结果与日志文件的目录。包含以下子目录：
    \begin{itemize}
      \item \texttt{in-design.redhawk} 或 \texttt{in-design.redhawk\_sc}：其中含 \texttt{design\_name.result}，RedHawk 或 RedHawk-SC 在此保存分析结果文件。
      \item \texttt{design\_name}：RedHawk 或 RedHawk-SC 写出分析结果供 Fusion Compiler 检索以显示地图与查询数据的目录。例如，运行 \cmd{open\_rail\_result} 会加载该目录中的数据以便在 GUI 中显示地图。
    \end{itemize}
\end{itemize}

若分析完成且希望保留数据供稍后检索，需在进行另一次 \cmd{analyze\_rail} 运行之前，使用 \appopt{rail.database} 将现有轨数据保存到另一工作目录。例如：
\begin{lstlisting}
fc_shell> set_app_options -name rail.database \
   -value RAIL_DATABASE_STATIC_RUN
\end{lstlisting}

否则，RedHawk/RedHawk-SC 工具会用新数据覆盖先前生成的轨数据，并发出如下警告：
\begin{lstlisting}
Warning: Rail result still open! Will be overwritten!
\end{lstlisting}

要避免该警告，在另一次 \cmd{analyze\_rail} 之前运行 \cmd{close\_rail\_result}，从内存中移除所有轨数据。
